%%%%%%%%%%%%%%%%%%%%%%%%%%%%%%%%%%%%%%%%%%%%%%%%%%%%%%%%%%%%%%%%%%%%%%%%%%%%%%%
% APÉNDICE A - REGISTROS DEL LEVANTAMIENTO DE CAMPO
%%%%%%%%%%%%%%%%%%%%%%%%%%%%%%%%%%%%%%%%%%%%%%%%%%%%%%%%%%%%%%%%%%%%%%%%%%%%%%%

\chapter{Registros de la Balanza}
\label{ap:registros_campo}

\normalsize

\noindent
\textbf{Lugar:} Panadería San Miguel\\
\textbf{Jornada ordinaria:} 02-09-2026\\
\textbf{Jornada de sábado:} 05-09-2026

\vspace{0.5cm}

El presente apéndice reúne los registros obtenidos durante el
levantamiento de campo del procedimiento manual de dosificación de harina.
Los datos corresponden a una jornada ordinaria y a una jornada de sábado.

%%%%%%%%%%%%%%%%%%%%%%%%%%%%%%%%%%%%%%%%%%%%%%%%%%%%%%%%%%%%%%%%%%%%%%%%%%%%%%%
% LECTURAS ACUMULADAS DE LA BALANZA
%%%%%%%%%%%%%%%%%%%%%%%%%%%%%%%%%%%%%%%%%%%%%%%%%%%%%%%%%%%%%%%%%%%%%%%%%%%%%%%

\section{Lecturas acumuladas de la balanza}

Las lecturas registradas corresponden a la Harina acumulada indicada por
la balanza después de cada carga principal. La indicación inicial,
después de aplicar la función de tara, corresponde a
$0{,}000\ \mathrm{kg}$.

\begin{table}[H]
    \centering
    \caption{Lecturas acumuladas de la balanza durante los ciclos de dosificación}
    \label{tab:lecturas_acumuladas}
    \small
    \setlength{\tabcolsep}{3pt}
    \renewcommand{\arraystretch}{1.15}

    \begin{tabular}{lccccc}
        \hline
        \textbf{Lote} &
        \textbf{\shortstack{Tras\\carga 1}} &
        \textbf{\shortstack{Tras\\carga 2}} &
        \textbf{\shortstack{Tras\\carga 3}} &
        \textbf{\shortstack{Tras\\carga 4}} &
        \textbf{\shortstack{Lectura\\final}} \\
        \hline

        \multicolumn{6}{c}{\textbf{Jornada ordinaria --- 02-09-2026}} \\
        \hline

        D1-C1 & $2{,}180$ & $4{,}530$ & $6{,}325$ & --        & $6{,}520$ \\
        D1-M1 & $2{,}090$ & $4{,}220$ & $6{,}348$ & $7{,}842$ & $8{,}018$ \\
        D1-F1 & $2{,}160$ & $4{,}315$ & $6{,}438$ & $7{,}368$ & $7{,}516$ \\
        D1-C2 & $2{,}240$ & $4{,}465$ & $6{,}342$ & --        & $6{,}513$ \\
        D1-M2 & $2{,}215$ & $4{,}386$ & $6{,}512$ & $7{,}918$ & $8{,}011$ \\

        \hline
        \multicolumn{5}{r}{\textbf{Total jornada ordinaria}} &
        $\mathbf{36{,}578}$ \\
        \hline

        \multicolumn{6}{c}{\textbf{Jornada de sábado --- 05-09-2026}} \\
        \hline

        S-C1 & $2{,}188$ & $4{,}680$ & $6{,}330$ & -- & $6{,}538$ \\
        S-E1 & $2{,}140$ & $4{,}300$ & $4{,}910$ & -- & $5{,}027$ \\

        \hline
        \multicolumn{5}{r}{\textbf{Total sábado}} &
        $\mathbf{11{,}565}$ \\

        \multicolumn{5}{r}{\textbf{Total general}} &
        $\mathbf{48{,}143}$ \\
        \hline
    \end{tabular}

    \vspace{0.3em}

    {\footnotesize Fuente: Elaboración propia.}
\end{table}

El símbolo ``--'' indica que el ciclo correspondiente concluyó las cargas
principales sin realizar una cuarta carga.

La identificación de los lotes utiliza las siguientes abreviaturas:
D1 corresponde a la jornada ordinaria; S corresponde al sábado; C
corresponde a pan casero; M a marraqueta; F a pan francés; y E a la Harina
común destinada a empanadas y rollos de queso. El lote S-E1 corresponde
a un único ciclo de dosificación.

%%%%%%%%%%%%%%%%%%%%%%%%%%%%%%%%%%%%%%%%%%%%%%%%%%%%%%%%%%%%%%%%%%%%%%%%%%%%%%%
% CANTIDADES DE HARINA Y PRODUCCIÓN
%%%%%%%%%%%%%%%%%%%%%%%%%%%%%%%%%%%%%%%%%%%%%%%%%%%%%%%%%%%%%%%%%%%%%%%%%%%%%%%

\section{Cantidades registradas durante la jornada ordinaria}

\begin{table}[H]
    \centering
    \caption{Cantidades de harina y producción registradas durante la jornada ordinaria}
    \label{tab:cantidades_jornada_ordinaria}
    \small
    \setlength{\tabcolsep}{3pt}
    \renewcommand{\arraystretch}{1.15}

    \begin{tabular}{lcccccc}
        \hline

        \textbf{Producto} &
        \textbf{\shortstack{Harina de\\receta (kg)}} &
        \textbf{\shortstack{Unidades de\\receta}} &
        \textbf{\shortstack{Unidades\\producidas}} &
        \textbf{\shortstack{Tiempo\\total (s)}} &
        \textbf{\shortstack{Harina final\\(kg)}} &
        \textbf{\shortstack{Diferencia\\(kg)}} \\

        \hline

        Casero --- 2 ciclos &
        $13{,}000$ & 400 & 405 &
        $107{,}1$ & $13{,}033$ & $+0{,}033$ \\

        Marraqueta --- 2 ciclos &
        $16{,}000$ & 300 & 305 &
        $130{,}8$ & $16{,}029$ & $+0{,}029$ \\

        Francés --- 1 ciclo &
        $7{,}500$ & 132 & 135 &
        $66{,}6$ & $7{,}516$ & $+0{,}016$ \\

        \hline

        \textbf{Total} &
        $\mathbf{36{,}500}$ &
        \textbf{832} &
        \textbf{845} &
        $\mathbf{304{,}5}$ &
        $\mathbf{36{,}578}$ &
        $\mathbf{+0{,}078}$ \\

        \hline
    \end{tabular}

    \vspace{0.3em}

    {\footnotesize Fuente: Elaboración propia.}
\end{table}

%%%%%%%%%%%%%%%%%%%%%%%%%%%%%%%%%%%%%%%%%%%%%%%%%%%%%%%%%%%%%%%%%%%%%%%%%%%%%%%
% SÁBADO
%%%%%%%%%%%%%%%%%%%%%%%%%%%%%%%%%%%%%%%%%%%%%%%%%%%%%%%%%%%%%%%%%%%%%%%%%%%%%%%

\section{Cantidades registradas durante la jornada de sábado}

\begin{table}[H]
    \centering
    \caption{Cantidades de harina y producción registradas durante la jornada de sábado}
    \label{tab:cantidades_sabado}
    \small
    \setlength{\tabcolsep}{3pt}
    \renewcommand{\arraystretch}{1.15}

    \begin{tabular}{lcccccc}
        \hline

        \textbf{Producto } &
        \textbf{\shortstack{Harina de\\receta (kg)}} &
        \textbf{\shortstack{Unidades de\\receta}} &
        \textbf{\shortstack{Unidades\\producidas}} &
        \textbf{\shortstack{Tiempo\\total (s)}} &
        \textbf{\shortstack{Harina final\\(kg)}} &
        \textbf{\shortstack{Diferencia\\(kg)}} \\

        \hline

        Casero &
        $6{,}500$ &
        200 panes &
        202 panes &
        $54{,}2$ &
        $6{,}538$ &
        $+0{,}038$ \\

        Harina común &
        $5{,}000$ &
        \shortstack{100 empanadas\\8 rollos} &
        \shortstack{102 empanadas\\8 rollos} &
        $53{,}2$ &
        $5{,}027$ &
        $+0{,}027$ \\

        \hline

        \textbf{Total} &
        $\mathbf{11{,}500}$ &
        -- &
        -- &
        $\mathbf{107{,}4}$ &
        $\mathbf{11{,}565}$ &
        $\mathbf{+0{,}065}$ \\

        \hline
    \end{tabular}

    \vspace{0.3em}

    {\footnotesize Fuente: Elaboración propia.}
\end{table}

La Harina común registrada constituye un único ciclo de dosificación. Las
empanadas y los rollos de queso se elaboraron a partir de esta misma Harina,
por lo que no se contabilizan como ciclos independientes.